% Speaker script for slides_1005.tex (meeting 5 Oct 2026), English + Japanese.
% Build: .\build_slides.ps1 -Tex speaker_notes_1005.tex -MaxPages 30   (lualatex, luatexja)
\documentclass[11pt,a4paper]{article}
\usepackage{luatexja}
\usepackage[margin=22mm]{geometry}
\usepackage{amsmath,amssymb}
\usepackage{xcolor}
\usepackage{enumitem}
\definecolor{c3}{RGB}{30,80,160}
\setlength{\parindent}{0pt}
\setlength{\parskip}{4pt}
\newcounter{sl}
\newcommand{\slide}[2]{\stepcounter{sl}\vspace{8pt}%
  {\color{c3}\large\bfseries Slide \thesl\ --- #1}\hfill{\small #2}\par\vspace{2pt}}
\newcommand{\EN}{\textbf{EN}\quad}
\newcommand{\JA}{\par\textbf{JA}\quad}

\begin{document}
{\Large\bfseries Speaker script --- meeting 5 October 2026}\\[2pt]
Assoc.\ Prof.\ Muramatsu $\times$ Prof.\ Soleimani $\times$ Nishioka.
16 slides plus a 4-slide appendix (notation, background, Klempt et al.\ 2024 reproduction), about 22 minutes. Times are rough targets.

% ------------------------------------------------------------------
\slide{Title}{0:20}
\EN Thank you for making time. Today I will show where the thesis stands.
The biofilm growth model is now coupled into the existing IKM element,
following the papers by Klempt et al.
\JA お時間をいただきありがとうございます。今日は修論の現状を報告します。
バイオフィルムの成長モデルを IKM の既存の要素に組み込み、すべて Klempt らの論文に
沿って進めた結果です。

% ------------------------------------------------------------------
\slide{Since the last meeting: the scope is fixed}{1:30}
\EN The main change since the last meeting is that the scope of the FEM part
is now fixed. The thesis covers one and two species, coupled into Oliver's
element as Klempt et al.\ 2024 describe it, with the species
composition taken from the point model of Klempt et al.\ 2026, and all
parameters from the papers. More species, starting from the four-species
cases of the 2026 paper and going up to the five oral species, would be continued
at Keio. For one and two species the papers give numbers I can check
against; for more species they do not yet.
\JA 前回からの一番大きな変化は、FEM 部分の範囲が決まったことです。修論では
1 菌種と 2 菌種を扱い、Oliver の要素に Klempt 2024 のとおりに組み込みます。
菌種の組成は Klempt 2026 の点モデルから取り、パラメータはすべて論文の値です。
それ以上の菌種数、つまり 2026 年論文の 4 菌種から始めて口腔の 5 菌種までは、
慶應で続ける計画です。1・2 菌種なら論文に照合できる数値がありますが、
それ以上の菌種数にはまだありません。

% ------------------------------------------------------------------
\slide{Verification of the growth law}{1:15}
\EN Before coupling anything, the growth law itself was checked. Against closed
forms: the fully constrained cube matches exactly, and free growth gives zero
stress to $10^{-10}$. The multi-species point model reproduces all fourteen numerical
examples of the 2026 paper, where I also found two misprints, and it
reproduces the figures of Fritsch et al.\ within 0.013.
\JA 組み込む前に、成長則そのものを確認しました。閉じた解との比較では、
完全拘束の立方体は厳密に一致し、自由成長では応力が $10^{-10}$ でゼロになります。
多菌種の点モデルは 2026 年論文の数値例
14 個をすべて再現し、その過程で論文の誤植を 2 つ見つけました。Fritsch らの図とも
0.013 以内で一致します。

% ------------------------------------------------------------------
\slide{The coupling: what the element does and what is added}{2:00}
\EN This is the coupling. Oliver's element solves the biofilm field of
Klempt 2024, Equation 34: diffusion, the source from growth, and front
growth driven by the nutrient. I do not change it. What I add happens at
each Gauss point of the material routine: I take $\phi$ from the field,
integrate Equation 36, $\dot\alpha=k_\alpha\phi$, and turn $\alpha$ into the
growth tensor and the stress with the verified law. The rate $k_\alpha$ is the
published value, $10^{-3}$ per unit time, used both in the field and in
Equation 36.
The model is a 2 mm cube with 512 elements. Only rigid-body motion is
suppressed, at three corner nodes, so uniform growth gives no stress and the
stress comes from differences in growth. The nutrient is held at $c=1$ in the
element layer on one face, and the biofilm starts as a seed of 32 elements
at the centre.
\JA これが組み込みの全体です。Oliver の要素は Klempt 2024 の式 34、つまり拡散、
成長による生成、栄養で駆動される界面成長を解きます。ここには手を入れていません。
私が加えたのは材料ルーチンの各ガウス点での処理です。場から $\phi$ を受け取り、
式 36 $\dot\alpha=k_\alpha\phi$ を積分し、その $\alpha$ を検証済みの成長則で
成長テンソルと応力にします。$k_\alpha$ は論文の値 $10^{-3}$ で、場と式 36 の両方で同じ値です。
モデルは 2 mm の立方体、512 要素です。力学は 3 つの角の節点で剛体運動だけを止めているので、
一様に成長すれば応力はゼロで、応力は成長の差から生じます。栄養は $y=-1$ mm の面の要素層で
$c=1$ に固定し、バイオフィルムは中心の 32 要素を $\phi=1$ で始めます。

% ------------------------------------------------------------------
\slide{Model and boundary conditions}{0:45}
\EN Here is the model in 3D. The seed of 32 elements is at the centre, the
nutrient is held on the layer at the back face, and only the three marked
corner nodes are constrained, which removes the rigid-body motion and
nothing more.
\JA モデルの 3 次元図です。中心に 32 要素の種があり、栄養は奥の面の要素層で固定して
います。拘束は印を付けた角の 3 節点だけで、剛体運動を止める以外の拘束はありません。

% ------------------------------------------------------------------
\slide{What I implemented, and what I used}{1:30}
\EN To be clear about what is mine: the growth field and the stiffness are
Oliver's element, used as it is; I read it and reported what I found. The
equations of the point model are from the paper; the solver is my own JAX
code, which reproduces all 14 examples of the paper. Everything else on this
table I wrote: the growth and composition step inside the element's material
call, the bridge from Fortran to Python, the sub-stepping, the scripts that
build and run the custom ANSYS, and the checks. On the Python side, each
Gauss point sends one request per step; splitting the step into point-model
steps of at most $10^{-4}$ and running them as one compiled loop brought a
single-element run from about six minutes to under six seconds, and every
call can be replayed in Python bit for bit.
\JA 自分の担当をはっきりさせておきます。成長場と剛性は Oliver さんの要素をそのまま
使い、私は中身を読んで見つけた点を報告しました。点モデルの式は論文のもので、解く
プログラムは私の JAX のコードです。論文の 14 例をすべて再現しています。表のそれ以外は
私が書いたものです。要素の材料ルーチンの中の成長と組成の処理、Fortran から Python への
橋渡し、時間の細分化、ANSYS のビルドと実行のスクリプト、そして検証です。Python 側では、
各積分点が 1 ステップに 1 回だけ要求を送ります。1 ステップを $10^{-4}$ 以下に分けて
まとめてコンパイルして回すことで、1 要素の計算が約 6 分から 6 秒弱になり、
すべての呼び出しは Python でビット単位まで再現できます。

% ------------------------------------------------------------------
\slide{Checks 1 and 2: exact solution and whole model}{2:00}
\EN Two checks of the coupling in ANSYS. On the left, an exact solution: where
the field has no gradients the equations reduce to a linear system with the
solution $\sinh$ and $\cosh$, and ANSYS follows it to round-off in $\alpha$
and to $10^{-7}$ in $\phi$, because this element still sees a little gradient.
On the right, the whole model: the stress computed with Equation 36 and with
the element's own growth variable has the same pattern, cosine similarity
1.0000, and the element's variable gives 0.455 times the result. That is one
half times ten elevenths: it averages over both species fields, including the
unused one, and it lags one time step.
\JA ANSYS での連成の確認を 2 つ示します。左は厳密解です。場に勾配がない所では式が
線形系になり、解は sinh と cosh です。ANSYS はこれに $\alpha$ で丸め誤差まで、
$\phi$ で $10^{-7}$ まで一致します。この要素も周りからわずかに勾配を受けているためです。
右はモデル全体です。式 36 で計算した応力と要素独自の成長変数で計算した応力は、
分布の形が同じでコサイン類似度は 1.0000、大きさは要素の変数のほうが 0.455 倍です。
これは 1/2 × 10/11 で、使っていない菌種の場も含めて平均していることと、
1 ステップ遅れていることで説明できます。

% ------------------------------------------------------------------
\slide{Check 3: the seeded element}{1:30}
\EN Third, the seeded element over time. Its growth $\alpha-1$ increases linearly to
$1.1\times10^{-3}$, which is exactly $k_\alpha\phi t$. The element itself is
compressed, and its neighbours carry about 25 times its von Mises stress.
Klempt et al.\ report pressure inside the biofilm and a ring of tension at its
edge in their 2024 paper; the compressed seed agrees with that, and whether the
neighbours are in tension I still have to read from their mean stress. The stresses are tiny, at most $10^{-4}$ Pa, because the published
growth rate is slow and the biofilm is very soft, $E=10$ Pa, as in their
Table 2.
\JA 3 つ目は、種を置いた要素の時間変化です。成長量 $\alpha-1$ は直線的に増えて
$1.1\times10^{-3}$ になり、これはちょうど $k_\alpha\phi t$ です。要素自体は圧縮され、
隣の要素はその約 25 倍の von Mises 応力を持ちます。Klempt らの 2024 年の論文では、
バイオフィルムの内部は圧縮、縁は引張の輪になっています。種の要素が圧縮されている点は
それと一致し、隣の要素が引張かどうかは平均応力でまだ確認していません。応力は最大 $10^{-4}$ Pa と非常に
小さいですが、これは論文の成長速度が遅く、バイオフィルムがとても柔らかい
（Table 2 の $E=10$ Pa）ためです。

% ------------------------------------------------------------------
\slide{Four things found along the way}{1:45}
\EN Four findings. First, the element's growth variable is the mean of the two
species fields, so with one species growth is halved. I would like to ask
Oliver whether that is intended. Second, start values must satisfy
$\phi_1+\phi_2\le1$, because the void fraction is one minus all
species; the example input starts both species at one in the seed, so for one
species I start the second at zero. Third, the first result set is not in
equilibrium, because the element returns zero stress while it sets up its
operators, so I exclude it. Fourth, the element's front-growth term is
inactive. When the input is read, the weights of the other species are written
into the weight of species 1, and the last one is zero, so this term is always
zero. In addition, it multiplies the Laplacian of $\phi$, where Equation 34 has
the gradient magnitude. In a test copy with both corrected, the seed spreads
as in the paper's example, but $\phi$ goes above 1. The three checks do not
depend on this term. I would like to discuss both points with Oliver before
using the spreading in the thesis. Earlier bring-up runs used larger growth rates
and the example input's stiffness, 1000 MPa; everything shown today uses the
published values.
\JA 見つかったことが 4 つあります。1 つ目、要素の成長変数は 2 つの菌種の場の平均なので、
1 菌種だと成長が半分になります。意図したものか Oliver に確認したいと思います。
2 つ目、空隙率は 1 から全菌種を引いたものなので、初期値は $\phi_1+\phi_2\le1$
でなければなりません。例題の入力では種の領域で両菌種とも 1 から始まっているので、
1 菌種の計算では 2 つ目を 0 にしました。3 つ目、最初の結果は釣り合っていません。
要素が演算子を準備している間は応力ゼロを返すためで、この結果は除いています。
4 つ目、要素の界面成長の項が働いていません。入力を読み込むとき、ほかの菌種の重みが
菌種 1 の重みに書き込まれていて、最後に読む値が 0 なので、この項は常に 0 になります。
さらに、この項は $\phi$ のラプラシアンを掛けていますが、式 34 では勾配の大きさです。
両方を直した試験用のコピーでは、論文の例と同じように種が広がりましたが、$\phi$ が 1 を
超えました。3 つの確認はこの項に依存しません。広がりを修論で使う前に、この 2 点を
Oliver と相談したいと思います。
なお、立ち上げ段階では大きな成長速度と例題入力の剛性（1000 MPa）を使っていましたが、
今日お見せした結果はすべて論文の値です。

% ------------------------------------------------------------------
\slide{Two species: composition from the point model}{1:45}
\EN For two species, the question is how to connect the point model to the 3-D
field. I chose a split that keeps both papers as they are: the field decides the
amount, $\phi_{3D}$, and drives growth; the point model decides only the
composition, that is which species and $\psi$. At each step its state is rescaled
so that the species add up to the amount from the field. The composition does
not act back on the stress, because in Klempt 2024 no material parameter
depends on the species. On ANSYS, every point and every step equals the scheme
run on its own, with difference zero.
\JA 2 菌種では、点モデルと 3 次元の場をどうつなぐかが問題です。両方の論文に忠実な
つなぎ方として、量 $\phi_{3D}$ は場が決めて成長を駆動し、点モデルは組成、つまり
どの菌種か、そして $\psi$ だけを決めます。各ステップで点モデルの状態を、
菌種の和が場の量に一致するように合わせ直します。組成は応力に戻りません。
Klempt 2024 では菌種によって変わる材料定数がないからです。ANSYS 上では、
全点・全ステップが単独で回した同じ計算と差 0 で一致しました。

% ------------------------------------------------------------------
\slide{Two species: results in the packed seed}{1:30}
\EN Results in the fully packed seed, for two cases of Table 1 with the paper's
parameters. Case 3 is coexistence: the paper gives 0.60 for species 1; coupled
I get 0.667, and between 0.664 and 0.668 when the cap is varied. Case 6 is
the case where one species wins: the paper has species 1 at zero with
$\psi_1=0.05$; I get 0.021 and $\psi_1=0.048$. So the type of outcome is
kept; the ratio shifts because the amount now comes from the 3-D field, not
from the point model filling up on its own. The result is converged in the
coupling time step.
\JA 種の領域、つまり完全に詰まった所での結果です。Table 1 の 2 ケースを論文の
パラメータで計算しました。ケース 3 は共存で、論文では菌種 1 が 0.60、組み込んだ計算では
0.667、上限値を変えても 0.664〜0.668 です。ケース 6 は片方が勝つケースで、論文では
菌種 1 がゼロ、$\psi_1=0.05$、こちらは 0.021 と 0.048 です。結果の種類は保たれていて、
比が変わるのは、量が点モデル自身の増え方ではなく 3 次元の場から来るためです。
結合の時間刻みに対して収束しています。

% ------------------------------------------------------------------
\slide{Two modelling assumptions}{1:45}
\EN Two assumptions are needed, and I state them explicitly. First, time: the
2024 paper runs to $T^*=1$, the 2026 paper counts steps of $10^{-4}$ in its own
unit, and no paper links the two. So the point model advances $s$ times the
time step, with $s=0.15$. In case 6 the winner only emerges for $s$ around
0.15 or larger; in case 3 it moves the result by 0.03 to 0.06. Second, in
fully packed regions the point model sees at most 0.9 of the amount; without
that, the result depends on the coupling time step. Both come with a
sensitivity study, and all other parameters are the published values.
\JA 仮定が 2 つ必要で、それを明示しています。1 つ目は時間です。2024 年論文は
$T^*=1$ まで、2026 年論文は独自の単位で $10^{-4}$ のステップを数えていて、
両者を結ぶ記述はどの論文にもありません。そこで点モデルは時間刻みの $s$ 倍だけ進むとし、
$s=0.15$ としました。これは結果に効きます。ケース 6 では $s$ が 0.15 程度以上でないと
勝者が現れず、ケース 3 では結果が 0.03〜0.06 動きます。2 つ目は、完全に詰まった所では
点モデルが見る量を最大 0.9 にすることです。これがないと結果が結合の時間刻みに依存します。
どちらも感度を調べてあり、それ以外のパラメータはすべて論文の値です。

% ------------------------------------------------------------------
\slide{Limits of the parts I added}{1:45}
\EN These are the limits I know of in my own part. First, the composition
is local: a point that the biofilm reaches starts from the seed composition,
it is not carried with the front. I tested how much this matters by
starting such points from their neighbour's state instead: the mean
composition changes by at most 0.001, because the diffuse edge of the field
reaches each point early. Second, the time link is assumed. Third, the
coupling inherits the behaviour of Eq.\ 34 as printed. Fourth, two growth
variables are integrated side by side; Check 2 showed how they relate.
\JA 自分の追加部分について、分かっている限界です。1 つ目、組成は各点で独立していて、
バイオフィルムが新しく届いた点は初期の組成から始まります。前線と一緒に運ばれません。
隣の点の状態から始めた場合と比べたところ、平均の組成の差は 0.001 以下でした。
場の薄い縁が早い時点で各点に届くためです。2 つ目、時間の対応は仮定です。3 つ目、
連成は印刷どおりの Eq.\ 34 の振る舞いをそのまま受け継ぎます。4 つ目、成長の変数が
2 つ並行して積分されています。その関係は Check 2 で示しました。

% ------------------------------------------------------------------
\slide{(Appendix D, only if asked) Klempt et al.\ 2024 test cases, reproduced}{1:45}
\EN I also reproduced the two test cases of Klempt et al.\ 2024 in Python.
With Table 2 unchanged, all three curves come within 0.02 to 0.09, under
two departures from the printed equations and one time scale per
simulation, which the paper's definition of $T^*$ allows. Three points I
would like to ask about: which form of Eq.\ 34 and 35 produced the figures,
the reference time of each run, and the unit of $\mu$. The pressure in
Table 3 fits $\mu$ read as MPa, not Pa.
\JA Klempt 2024 の 2 つのテストケースも Python で再現しました。Table 2 の値を変えずに、
3 本の曲線がすべて 0.03〜0.09 で合います。ただし、印刷された式から 2 か所変えています。
また、シミュレーションごとに時間スケールを決めていますが、これは論文の $T^*$ の定義で
許されています。お聞きしたいのは 3 点です。図を作った Eq.\ 34・35 の形、各計算の基準時間、
そして $\mu$ の単位です。Table 3 の圧力は、$\mu$ を Pa ではなく MPa として読むと合います。

% ------------------------------------------------------------------
\slide{Status and timeline}{1:00}
\EN On the administrative side, the admission was granted on 18 August, with
Prof.\ Junker and Dr.\ Wangenheim as examiners, and the internship at NIFE is
finished. The coupling for one and two species is done. October is for the
thesis runs and figures, writing runs in parallel, and I submit within
November, aiming for the beginning of the month. The presentation follows in
December, before I return to Japan; I would like to fix the date with both
examiners soon.
\JA 事務的には、8 月 18 日に登録が認められ、審査は Junker 先生と Wangenheim 先生です。
NIFE でのインターンも終わっています。1・2 菌種の組み込みは終わっていて、10 月で修論用の
計算と図を仕上げ、並行して執筆し、11 月中に提出します（目標は上旬）。発表は帰国前の
12 月で、日程を早めに両先生と決めたいと考えています。

% ------------------------------------------------------------------
\slide{Questions for this meeting}{2:00}
\EN I have six questions. First, the scope: is it agreed that the thesis
covers one and two species as published, and more species continue
at Keio? Second, are the two assumptions, the time link and the cap, each
with a sensitivity study, acceptable for the thesis? Third, is anything
needed from the Keio side for the double-degree requirements this term?
Fourth, on Klempt et al.\ 2024: the form of Eq.\ 34 and 35 behind the
figures, the reference time of each run, whether $\mu$ is in Pa or MPa, and
whether $\alpha$ follows Eq.\ 36 or the update in Table~1. And fifth, the
structure of the thesis: the TMCMC calibration of the point model, a short
section on the metabolic sign prior from my internship, then the ANSYS
coupling. Is it acceptable to include the internship part? And sixth, the
submission: how many printed copies, whether an electronic version is needed,
and whether anything has to be submitted on the Keio side.
\JA 6 つお聞きしたいことがあります。1 つ目、範囲について、修論は論文どおりの 1・2 菌種、
それ以上は慶應で続けるという分け方でよろしいでしょうか。2 つ目、時間の対応と上限値という
2 つの仮定を、感度の検討付きで修論に入れてよいでしょうか。3 つ目、ダブルディグリーの
要件で、今学期に慶應側で必要なことはありますか。4 つ目、Klempt 2024 について、
図を作った Eq.\ 34・35 の形、各計算の基準時間、$\mu$ が Pa か MPa か、$\alpha$ は
Eq.\ 36 と Table~1 のどちらに従うのかを教えてください。5 つ目、修論の構成です。
TMCMC による点モデルの推定、インターンでの代謝からの符号の事前分布の短い節、
そのあと ANSYS への組み込み、という順にしたいと考えています。インターンの部分を
入れてもよいでしょうか。6 つ目、提出について、製本の部数、電子版の要否、
慶應側にも提出が必要なものがあるかを教えてください。

% ------------------------------------------------------------------
\slide{References}{(no talk)}
\EN Leave this slide up only if asked; the four papers are the ones cited on
the slides. Three appendix slides follow (notation, background of the growth
field, background of the point model and one coupling step), for questions only.
\JA 聞かれたときだけ表示します。スライドで引用した 4 本の論文です。その後に付録が
3 枚（記号の一覧、成長場の背景、点モデルの背景と 1 ステップの流れ）あり、
質問されたときに使います。

% ------------------------------------------------------------------
\vspace{12pt}
{\color{c3}\large\bfseries Possible questions and short answers}

\begin{description}[leftmargin=0pt,itemsep=6pt]
\item[Why are the stresses so small?]
\EN The published $k_\alpha=10^{-3}$ with $T^*\le1$ gives $\alpha-1\le10^{-3}$, and the
published stiffness is $E=10$ Pa, weighted with $\phi^2$. Stress scales
with both, so other values give other magnitudes, but they would no longer be
the paper's model. (I checked $\nu=0.45$ against 0.49: the stress changes by
27\,\% while the bulk modulus changes fivefold, so no sign of locking.)
\JA 論文の $k_\alpha=10^{-3}$ と $T^*\le1$ では $\alpha-1\le10^{-3}$ にしかならず、剛性も
論文では $E=10$ Pa で $\phi^2$ に比例します。応力はどちらにも比例するので、
値を変えれば大きさも変わりますが、それはもう論文のモデルではありません。
（$\nu=0.45$ と 0.49 を比べると、体積弾性率は 5 倍変わるのに応力の変化は 27\,\% なので、
要素が硬くなりすぎる兆候はありません。）

\item[Why not the five species right away?]
\EN There is no published reference beyond four species, and Oliver's element is
written for two fields; extending it is a rewrite plus new verification.
\JA 4 菌種を超えると論文の参照値がなく、Oliver の要素も 2 つの場を前提に書かれています。
拡張は書き直しと新たな検証が必要です。

\item[Is the composition result physical, given it differs from the paper?]
\EN The type of outcome matches; the ratio differs because the amount is set
by the field. The point model alone reproduces the paper exactly, so the shift
is the coupling's effect, and it is converged and insensitive to the cap.
\JA 結果の種類は一致しています。比の違いは量を場が決めるためです。点モデル単独では
論文を厳密に再現しているので、この差は結合による効果で、収束しており上限値にも鈍感です。

\item[Why $s=0.15$?]
\EN It maps the 2026 paper's 1500 steps of $10^{-4}$ onto the 2024 paper's
horizon $T^*=1$. It is an assumption, which is why it is reported with a sweep.
\JA 2026 年論文の $10^{-4}$ × 1500 ステップを、2024 年論文の $T^*=1$ に対応させる値です。
仮定なので、感度の結果と一緒に示しています。
\end{description}
\end{document}
